% BEATRIX: Antworten auf Fasnacht-Fragebogen (Version 2)
% Überarbeitet nach Style-Guide: SSBM-Wording + Fehr-Safe Language
% Für Innosuisse Innovation Cheque (CHF 15'000)
% Datum: 2026-01-18

\documentclass[11pt,a4paper]{article}

% Packages
\usepackage[utf8]{inputenc}
\usepackage[T1]{fontenc}
\usepackage[ngerman]{babel}
\usepackage{geometry}
\usepackage{graphicx}
\usepackage{booktabs}
\usepackage{array}
\usepackage{longtable}
\usepackage{tabularx}
\usepackage{xcolor}
\usepackage{hyperref}
\usepackage{enumitem}
\usepackage{fancyhdr}
\usepackage{titlesec}
\usepackage{tcolorbox}
\usepackage{amsmath}
\usepackage{amssymb}

% Page geometry
\geometry{left=2.5cm, right=2.5cm, top=2.5cm, bottom=2.5cm}

% Colors - Innosuisse corporate
\definecolor{innosuisse}{RGB}{0, 102, 153}
\definecolor{lightgray}{RGB}{245, 245, 245}
\definecolor{checkgreen}{RGB}{0, 128, 0}
\definecolor{warnorange}{RGB}{230, 126, 34}

% Hyperref setup
\hypersetup{
    colorlinks=true,
    linkcolor=innosuisse,
    urlcolor=innosuisse,
    citecolor=innosuisse
}

% Header and footer
\pagestyle{fancy}
\fancyhf{}
\fancyhead[L]{\small BEATRIX -- Innovation Cheque}
\fancyhead[R]{\small Innosuisse SSBM}
\fancyfoot[C]{\thepage}
\renewcommand{\headrulewidth}{0.4pt}

% Section formatting
\titleformat{\section}{\Large\bfseries\color{innosuisse}}{\thesection}{1em}{}
\titleformat{\subsection}{\large\bfseries}{\thesubsection}{1em}{}
\titleformat{\subsubsection}{\normalsize\bfseries}{\thesubsubsection}{1em}{}

% Custom box for key insights
\newtcolorbox{insightbox}{
    colback=white,
    colframe=innosuisse,
    boxrule=1pt,
    arc=3pt,
    left=8pt,
    right=8pt,
    top=6pt,
    bottom=6pt
}

% Custom box for summaries
\newtcolorbox{summarybox}[1][]{
    colback=lightgray,
    colframe=innosuisse,
    fonttitle=\bfseries,
    title=#1,
    boxrule=0.5pt,
    arc=3pt
}

% Custom box for honest limitations
\newtcolorbox{limitbox}{
    colback=white,
    colframe=warnorange,
    boxrule=1pt,
    arc=3pt,
    left=8pt,
    right=8pt
}

\begin{document}

% Title page
\begin{titlepage}
    \centering
    \vspace*{1.5cm}

    {\Huge\bfseries\color{innosuisse} BEATRIX\par}
    \vspace{0.5cm}
    {\Large Behavioral Economics AI Technology for\\Research and Implementation eXcellence\par}

    \vspace{1.5cm}

    {\LARGE\bfseries Antworten auf den\\Innosuisse Innovation Cheque Fragebogen\par}

    \vspace{0.5cm}
    {\large\textit{Version 2.0 -- Überarbeitet nach Style-Guide}\par}

    \vspace{1.5cm}

    \begin{tabular}{ll}
        \textbf{Förderprogramm:} & Innosuisse Innovation Cheque \\
        \textbf{Fördersumme:} & CHF 15'000 \\
        \textbf{Fachgebiet:} & Social Sciences \& Business Management (SSBM) \\
        \textbf{Wirtschaftspartner:} & FehrAdvice Partners AG \\
        \textbf{Forschungspartner:} & Universität Zürich \\
    \end{tabular}

    \vspace{1.5cm}

    \begin{tabular}{ll}
        \textbf{Prof. Harald Gall} & Institut für Informatik, UZH \\
        \textbf{Prof. Johannes Luger} & Institut für Betriebswirtschaft, UZH \\
        \textbf{Prof. Ernst Fehr} & Scientific Advisor \\
    \end{tabular}

    \vfill

    \begin{tcolorbox}[colback=lightgray, colframe=innosuisse, width=0.9\textwidth]
    \centering
    \textbf{Strategischer Kernsatz}\\[0.3em]
    \textit{``Die Technologie ermöglicht die Innovation, aber die Transformation der Beratungspraxis ist das Produkt.''}
    \end{tcolorbox}

    \vspace{1cm}

    {\large Datum: 18. Januar 2026\par}
\end{titlepage}

% Table of contents
\tableofcontents
\newpage

%==============================================================================
% ABSTRACT
%==============================================================================
\section*{Abstract}
\addcontentsline{toc}{section}{Abstract}

Unternehmen und öffentliche Verwaltungen verfügen über exzellente Datensysteme für harte Faktoren (Finanzen, Logistik, Produktion), doch für die systematische Analyse von Verhaltenstreibern fehlen wissenschaftlich fundierte Werkzeuge. Die Folge: Etwa 70\% aller Change-Projekte erreichen ihre Ziele nicht, weil menschliche Faktoren unterschätzt werden.

BEATRIX ist ein evidenzbasiertes Entscheidungsunterstützungssystem, das 50 Jahre verhaltensökonomische Forschung für die Managementpraxis zugänglich macht. Die Innovation besteht aus drei Elementen: (1) Systematische Kontextmodellierung, (2) Analyse von Wechselwirkungen zwischen Massnahmen, und (3) transparente, nachvollziehbare Entscheidungslogik.

Die Universität Zürich validiert mit zwei komplementären Kompetenzen: Prof.~Harald Gall analysiert die technische Architektur auf Skalierbarkeit. Prof.~Johannes Luger untersucht die Nutzerakzeptanz und entwickelt Einführungsstrategien.

Der Innovation Cheque finanziert eine fokussierte Machbarkeitsstudie (125 Arbeitsstunden): ein Architektur-Assessment, eine Usability-Studie mit 10--15 Führungskräften, sowie einen Validierungsbericht als Grundlage für ein Innosuisse-Hauptprojekt.

BEATRIX adressiert einen relevanten Markt im DACH-Raum. Pilotprojekte zeigen vielversprechende Resultate. Das Projekt stärkt die Schweizer Position als Standort für evidenzbasierte Management-Innovationen.

\vspace{0.5cm}
\noindent\textit{Wortanzahl: 198}

\newpage

%==============================================================================
% FRAGE A: INNOVATION POTENTIAL
%==============================================================================
\section{Innovation Potential}

\subsection{Das Problem, das wir lösen}

Stellen Sie sich vor: Ein Unternehmen will die Mitarbeiterzufriedenheit verbessern. Der Berater empfiehlt zehn Massnahmen -- aber welche davon wirken zusammen? Welche blockieren sich gegenseitig? Und funktioniert das, was in Skandinavien klappt, auch in der Schweiz?

\begin{insightbox}
\textbf{Heute basiert diese Einschätzung auf Erfahrung und Intuition.} BEATRIX macht sie systematisch und nachvollziehbar.
\end{insightbox}

\subsection{Die drei Säulen der Innovation}

\subsubsection{Säule 1: Kontext ernst nehmen}

Menschen verhalten sich nicht im luftleeren Raum. Ein finanzieller Anreiz wirkt anders bei einem Berufseinsteiger als bei einer Führungskraft kurz vor der Pensionierung. BEATRIX berücksichtigt \textbf{acht Kontext-Dimensionen} systematisch:

\begin{table}[h]
\centering
\begin{tabularx}{\textwidth}{lX}
\toprule
\textbf{Dimension} & \textbf{Frage} \\
\midrule
Zeitlicher Kontext & Ist die Entscheidung dringend oder langfristig? \\
Sozialer Kontext & Wer beobachtet? Wer urteilt? \\
Institutioneller Kontext & Welche Regeln gelten? Welche Kultur herrscht? \\
Informationeller Kontext & Was weiss die Person? Was fehlt? \\
Emotionaler Kontext & Ist die Person gestresst? Entspannt? \\
Ressourcen-Kontext & Hat sie Zeit, Geld, Energie für die Entscheidung? \\
Kultureller Kontext & Schweiz ist nicht Schweden ist nicht Singapur \\
Physischer Kontext & Wo findet die Entscheidung statt? \\
\bottomrule
\end{tabularx}
\end{table}

\textbf{Warum ist das neu?} Bisherige Ansätze behandeln Kontext als Störfaktor. Wir machen ihn zum zentralen Ordnungsprinzip.

\subsubsection{Säule 2: Wechselwirkungen sichtbar machen}

Wenn ein Unternehmen gleichzeitig Boni einführt und auf Teamarbeit setzt, können sich diese Massnahmen verstärken -- oder gegenseitig aufheben.

\begin{table}[h]
\centering
\begin{tabularx}{\textwidth}{llX}
\toprule
\textbf{Kombination} & \textbf{Effekt} & \textbf{Beispiel} \\
\midrule
Verstärkend & Wirken zusammen stärker & Transparenz + Feedbackkultur \\
Neutralisierend & Heben sich auf & Individualboni + Teamziele \\
Blockierend & Eine verhindert die andere & Kontrolle + Vertrauenskultur \\
\bottomrule
\end{tabularx}
\end{table}

\textbf{Warum ist das neu?} Traditionelle Beratung betrachtet Massnahmen einzeln. Wir zeigen systematisch, welche Kombinationen in welchem Kontext wirken.

\subsubsection{Säule 3: Transparente Entscheidungslogik}

BEATRIX ist kein ``Black Box''-System. Jede Empfehlung ist nachvollziehbar und basiert auf dokumentierter Forschung.

\begin{table}[h]
\centering
\begin{tabular}{ll}
\toprule
\textbf{Traditionell} & \textbf{Mit BEATRIX} \\
\midrule
``Der Berater sagt, wir sollen X tun.'' & ``Die Analyse zeigt: In unserem Kontext wirkt X, weil Y.'' \\
Vertrauen in Person & Vertrauen in Methode \\
Schwer replizierbar & Nachvollziehbar und lernend \\
\bottomrule
\end{tabular}
\end{table}

\subsection{Was BEATRIX nicht ist}

\begin{limitbox}
Wir sind ehrlich über die Grenzen:
\begin{itemize}[noitemsep]
    \item \textbf{Kein Wahrsager:} BEATRIX macht keine Vorhersagen über individuelle Entscheidungen
    \item \textbf{Kein Ersatz für Führung:} Das System unterstützt Entscheidungen, es trifft sie nicht
    \item \textbf{Kein Allheilmittel:} Manche Probleme sind nicht durch bessere Analyse lösbar
    \item \textbf{Keine ``künstliche Intelligenz'' im Hype-Sinn:} Systematische Analyse, kein magisches Orakel
\end{itemize}
\end{limitbox}

\subsection{Was validiert der Forschungspartner?}

\begin{table}[h]
\centering
\begin{tabularx}{\textwidth}{lXl}
\toprule
\textbf{Wer} & \textbf{Fragt} & \textbf{Methode} \\
\midrule
\textbf{Prof. Gall} & Ist das System robust und skalierbar? & Technische Analyse \\
\textbf{Prof. Luger} & Vertrauen Manager den Empfehlungen? & Tests mit Führungskräften \\
\bottomrule
\end{tabularx}
\end{table}

\begin{summarybox}[Zusammenfassung: Innovation Potential]
\textbf{BEATRIX ist eine Management-Innovation, keine Software-Innovation.}

Das System weiterentwickelt die Beratungspraxis systematisch, indem es:
\begin{enumerate}[noitemsep]
    \item \textbf{Kontext} als zentrales Ordnungsprinzip verwendet
    \item \textbf{Wechselwirkungen} zwischen Massnahmen sichtbar macht
    \item \textbf{Transparente Entscheidungslogik} bietet statt ``Black Box''
\end{enumerate}

Die UZH prüft: Ist das System technisch robust (Gall)? Vertrauen Manager den Empfehlungen (Luger)?
\end{summarybox}

\newpage

%==============================================================================
% FRAGE B: MARKET RELEVANCE
%==============================================================================
\section{Market Relevance}

\subsection{Die Soft-Factor-Lücke}

Moderne Unternehmen haben für fast alles Daten:

\begin{table}[h]
\centering
\begin{tabular}{lll}
\toprule
\textbf{Bereich} & \textbf{Datenqualität} & \textbf{Werkzeuge} \\
\midrule
Finanzen & Exzellent & SAP, Oracle \\
Logistik & Sehr gut & Supply-Chain-Systeme \\
Marketing & Gut & CRM, Analytics \\
\textbf{Mitarbeiterverhalten} & \textbf{Lückenhaft} & \textbf{Umfragen, Intuition} \\
\textbf{Kundenentscheidungen} & \textbf{Lückenhaft} & \textbf{Fokusgruppen, Bauchgefühl} \\
\bottomrule
\end{tabular}
\end{table}

\begin{insightbox}
\textbf{Das Problem:} Für die ``harten'' Faktoren gibt es hervorragende Systeme. Für die ``weichen'' Faktoren -- menschliches Verhalten, Motivation, Entscheidungen -- fehlen systematische Werkzeuge.
\end{insightbox}

\subsection{Die Konsequenz}

Studien zeigen: \textbf{Etwa 70\% aller Change-Projekte erreichen ihre Ziele nicht} (McKinsey, Kotter). Der häufigste Grund? Menschliche Faktoren wurden unterschätzt.

\begin{table}[h]
\centering
\begin{tabularx}{\textwidth}{lXX}
\toprule
\textbf{Typisches Szenario} & \textbf{Was passiert} & \textbf{Mit systematischer Analyse} \\
\midrule
Neue Software einführen & Mitarbeiter nutzen sie nicht & Widerstände früh erkennen \\
Kultur verändern & Alte Gewohnheiten bleiben & Hebel identifizieren \\
Anreizsystem ändern & Unerwartete Nebenwirkungen & Wechselwirkungen prüfen \\
\bottomrule
\end{tabularx}
\end{table}

\subsection{Für wen ist BEATRIX gedacht?}

\begin{table}[h]
\centering
\begin{tabularx}{\textwidth}{lXX}
\toprule
\textbf{Wer?} & \textbf{Typische Frage} & \textbf{Wie BEATRIX hilft} \\
\midrule
Beratungsunternehmen & ``Wie können wir unsere Empfehlungen besser begründen?'' & Systematische Analyse statt Intuition \\
Personalabteilungen & ``Warum funktioniert unsere neue Strategie nicht?'' & Kontext-Analyse der Belegschaft \\
Öffentliche Verwaltung & ``Wie erreichen wir, dass Bürger das neue Angebot nutzen?'' & Evidenzbasierte Gestaltung \\
Geschäftsführung & ``Welche Massnahmen wirken wirklich?'' & Priorisierung nach Wirksamkeit \\
\bottomrule
\end{tabularx}
\end{table}

\subsection{Ein konkretes Beispiel}

\textbf{Situation:} Ein Energieversorger will Kunden zum Wechsel auf einen günstigeren Tarif bewegen.

\textbf{Traditioneller Ansatz:} Brief schicken, hoffen, dass Kunden reagieren.

\textbf{Mit systematischer Analyse:}
\begin{itemize}[noitemsep]
    \item Kontext verstehen: Warum wechseln Kunden nicht, obwohl sie sparen würden?
    \item Hindernisse identifizieren: Komplexität? Trägheit? Misstrauen?
    \item Massnahmen anpassen: Einfacherer Prozess? Persönliche Ansprache?
\end{itemize}

\textbf{Ergebnis in Pilotprojekten:} Deutlich höhere Wechselquoten -- weil die Hindernisse verstanden wurden.

\subsection{Marktgrösse}

\begin{table}[h]
\centering
\begin{tabularx}{\textwidth}{lXr}
\toprule
\textbf{Ebene} & \textbf{Beschreibung} & \textbf{Schätzung} \\
\midrule
Gesamtmarkt & Beratungsmarkt weltweit & USD 300+ Mrd. \\
Relevanter Markt & Organisations-/Verhaltensberatung (DACH) & CHF 5--10 Mrd. \\
Erreichbarer Markt & Was wir in 5 Jahren realistisch erreichen können & CHF 50--100 Mio. \\
\bottomrule
\end{tabularx}
\end{table}

\textit{Hinweis: Diese Schätzungen sind konservativ. Der tatsächliche Markt könnte grösser sein, aber wir vermeiden übertriebene Prognosen.}

\begin{summarybox}[Zusammenfassung: Market Relevance]
\textbf{BEATRIX adressiert eine echte Lücke:} Während für harte Faktoren exzellente Werkzeuge existieren, fehlt für Verhaltensfragen ein systematischer Ansatz.

\begin{itemize}[noitemsep]
    \item \textbf{Klares Problem:} Hohe Scheiterquote bei Veränderungsprojekten
    \item \textbf{Definierte Zielgruppen:} Beratungen, Personalabteilungen, öffentliche Verwaltung
    \item \textbf{Validierte Wirksamkeit:} Positive Ergebnisse in Pilotprojekten
    \item \textbf{Realistisches Marktpotenzial:} Konservativ geschätzt
\end{itemize}
\end{summarybox}

\newpage

%==============================================================================
% FRAGE C: FEASIBILITY & RISK REDUCTION
%==============================================================================
\section{Feasibility \& Risk Reduction}

\subsection{Was wurde bereits entwickelt?}

\begin{table}[h]
\centering
\begin{tabular}{llll}
\toprule
\textbf{Phase} & \textbf{Zeitraum} & \textbf{Was wurde erreicht?} \\
\midrule
Konzeption & 2018--2020 & Theoretisches Fundament erarbeitet \\
Prototyp & 2020--2022 & Erste funktionierende Version \\
Pilotphase & 2022--2025 & Einsatz bei über 20 Kundenprojekten \\
Heute & 2026 & Funktionierender Kern, bereit für Skalierung \\
\bottomrule
\end{tabular}
\end{table}

\begin{insightbox}
\textbf{Aktueller Stand:} Das System funktioniert und wird eingesetzt. Was fehlt, ist die externe Validierung der Skalierbarkeit und der Nutzerakzeptanz.
\end{insightbox}

\subsection{Was funktioniert bereits?}

\begin{itemize}
    \item[\checkmark] Systematische Erfassung der acht Kontext-Dimensionen
    \item[\checkmark] Modellierung von Wechselwirkungen zwischen Massnahmen
    \item[\checkmark] Benutzeroberfläche für Berater
    \item[\checkmark] Bibliothek mit über 50 dokumentierten Anwendungsfällen
    \item[\checkmark] Integration von über 1'500 wissenschaftlichen Studien
\end{itemize}

\subsection{Die zwei offenen Fragen}

\begin{table}[h]
\centering
\begin{tabularx}{\textwidth}{lXl}
\toprule
\textbf{Frage} & \textbf{Warum kritisch?} & \textbf{Wer prüft?} \\
\midrule
Ist das System skalierbar? & Ohne Skalierbarkeit kein Geschäftsmodell & Prof. Gall \\
Vertrauen Manager den Empfehlungen? & Ohne Akzeptanz keine Nutzung & Prof. Luger \\
\bottomrule
\end{tabularx}
\end{table}

\subsection{Was kann mit CHF 15'000 erreicht werden?}

Der Innovation Cheque finanziert \textbf{125 Arbeitsstunden}:

\begin{table}[h]
\centering
\begin{tabular}{lrl}
\toprule
\textbf{Arbeitspaket} & \textbf{Umfang} & \textbf{Ergebnis} \\
\midrule
Architektur-Analyse & 40 Stunden & Technischer Prüfbericht \\
Nutzer-Tests & 50 Stunden & Akzeptanz-Studie \\
Risikoanalyse & 20 Stunden & Identifikation kritischer Punkte \\
Abschlussbericht & 15 Stunden & Empfehlungen für Hauptprojekt \\
\midrule
\textbf{Total} & \textbf{125 Stunden} & \\
\bottomrule
\end{tabular}
\end{table}

\subsection{Entwicklungspfad}

\begin{center}
\begin{tabular}{rcl}
\textbf{Heute (2026)} & $\leftarrow$ & Wir sind hier \\
& $\downarrow$ & \\
\textbf{Innovation Cheque} & & Machbarkeit prüfen (CHF 15k) \\
& $\downarrow$ & \\
\textbf{Hauptprojekt} & & Skalierbare Lösung (CHF 500k--2M) \\
& $\downarrow$ & \\
\textbf{Marktreife (2028--29)} & & Breiter Einsatz möglich \\
\end{tabular}
\end{center}

\begin{summarybox}[Zusammenfassung: Feasibility \& Risk Reduction]
\textbf{BEATRIX ist keine Idee auf dem Papier} -- das System funktioniert und wird eingesetzt.

Der Innovation Cheque klärt zwei kritische Fragen:
\begin{itemize}[noitemsep]
    \item Ist das System technisch bereit für breiteren Einsatz? (Gall)
    \item Vertrauen Führungskräfte den Empfehlungen? (Luger)
\end{itemize}

Mit CHF 15'000 werden 125 Arbeitsstunden finanziert, die eine fundierte Go/No-Go-Entscheidung für das Hauptprojekt ermöglichen.
\end{summarybox}

\newpage

%==============================================================================
% FRAGE D: RESEARCH PARTNER COLLABORATION
%==============================================================================
\section{Collaboration with Research Partner}

\subsection{Was FehrAdvice mitbringt -- und was fehlt}

\begin{table}[h]
\centering
\begin{tabularx}{\textwidth}{lX}
\toprule
\textbf{Stärke von FehrAdvice} & \textbf{Was fehlt} \\
\midrule
Verhaltensökonomie (15+ Jahre) & Technische Architektur für Skalierung \\
Praxiswissen (20+ Projekte) & Wissenschaftliche Akzeptanzforschung \\
Kundenzugang & Forschungsbasiertes Change Management \\
Ernst Fehr als Gründer & Unabhängige externe Prüfung \\
\bottomrule
\end{tabularx}
\end{table}

\subsection{Warum die Universität Zürich?}

\begin{table}[h]
\centering
\begin{tabularx}{\textwidth}{lXc}
\toprule
\textbf{Kriterium} & \textbf{Anforderung} & \textbf{UZH} \\
\midrule
Wissenschaftliche Qualität & Höchste Standards & \checkmark~Weltklasse \\
Thematische Passung & Software + Management & \checkmark~Gall + Luger \\
Bestehende Beziehung & Vertrauensbasis & \checkmark~Fehr ist UZH-Professor \\
Nähe & Enge Zusammenarbeit möglich & \checkmark~Alle in Zürich \\
\bottomrule
\end{tabularx}
\end{table}

\subsection{Die zwei Professoren}

\textbf{Prof. Harald Gall} (Institut für Informatik)

Einer der weltweit führenden Experten für Software-Entwicklung. Seine Frage:
\begin{quote}
\textit{``Wie muss ein System wie BEATRIX gebaut sein, damit es zuverlässig wächst und sich weiterentwickelt?''}
\end{quote}

\textbf{Prof. Johannes Luger} (Institut für Betriebswirtschaft)

Forscht zur Zusammenarbeit zwischen Menschen und intelligenten Systemen. Seine Frage:
\begin{quote}
\textit{``Unter welchen Bedingungen vertrauen Führungskräfte den Empfehlungen eines Analyse-Systems?''}
\end{quote}

\subsection{Konkrete Beiträge}

\begin{table}[h]
\centering
\begin{tabularx}{\textwidth}{lXX}
\toprule
\textbf{Wer} & \textbf{Beitrag} & \textbf{Ergebnis} \\
\midrule
Prof. Gall & Architektur-Analyse & Prüfbericht, Skalierungs-Empfehlungen \\
Prof. Luger & Nutzer-Tests mit 10--15 Führungskräften & Akzeptanz-Studie, Einführungsstrategien \\
Beide & Risikoanalyse & Validierungsbericht für Hauptprojekt \\
\bottomrule
\end{tabularx}
\end{table}

\begin{summarybox}[Zusammenfassung: Research Partner Collaboration]
\textbf{Die UZH ist der ideale Partner} für dieses Projekt:

\begin{itemize}[noitemsep]
    \item \textbf{Ernst Fehr} verbindet FehrAdvice und UZH bereits
    \item \textbf{Prof. Gall} sichert die technische Qualität
    \item \textbf{Prof. Luger} erforscht die menschlichen Faktoren
    \item \textbf{Beide} bringen unabhängige wissenschaftliche Prüfung
\end{itemize}

Die Partnerschaft schliesst kritische Kompetenzlücken und legt den Grundstein für langfristige Zusammenarbeit.
\end{summarybox}

\newpage

%==============================================================================
% FRAGE E: IMPACT FOR SWITZERLAND
%==============================================================================
\section{Impact for Switzerland}

\subsection{Eine Schweizer Lösung}

\begin{table}[h]
\centering
\begin{tabularx}{\textwidth}{lX}
\toprule
\textbf{Aspekt} & \textbf{Bedeutung} \\
\midrule
Geistiges Eigentum bleibt in der Schweiz & Keine Abhängigkeit von ausländischen Anbietern \\
Daten unter Schweizer Recht & Vertrauenswürdigkeit für sensible Analysen \\
Lokale Wertschöpfung & Arbeitsplätze, Steuern, Know-how in der Schweiz \\
``Made in Switzerland'' & Qualitätsmerkmal auch für Beratungsdienstleistungen \\
\bottomrule
\end{tabularx}
\end{table}

\subsection{Arbeitsplätze und Wertschöpfung}

\begin{table}[h]
\centering
\begin{tabular}{lll}
\toprule
\textbf{Zeitraum} & \textbf{Erwartung} \\
\midrule
2026 (Machbarkeitsphase) & 2--3 direkte Arbeitsplätze \\
2028 (Markteinführung) & 10--15 direkte Arbeitsplätze \\
2030 (Skalierung) & 30--50 direkte Arbeitsplätze \\
\bottomrule
\end{tabular}
\end{table}

\textit{Hinweis: Diese Schätzungen sind konservativ und hängen vom Projekterfolg ab.}

\subsection{Die einzigartige Schweizer Position}

\begin{table}[h]
\centering
\begin{tabularx}{\textwidth}{lX}
\toprule
\textbf{Faktor} & \textbf{Schweiz-Vorteil} \\
\midrule
Ernst Fehr & Einer der weltweit führenden Verhaltensökonomen arbeitet hier \\
UZH-Forschung & Hervorragend positioniert in relevanten Feldern \\
FehrAdvice & Einziges Unternehmen dieser Art mit direkter Verbindung zur Spitzenforschung \\
Neutralität & Schweiz als vertrauenswürdiger Standort für sensible Analysen \\
Stabilität & Verlässliches Umfeld für langfristige Entwicklung \\
\bottomrule
\end{tabularx}
\end{table}

\subsection{Gesellschaftlicher Nutzen}

Wenn Organisationen systematischer mit Verhaltensfragen umgehen, profitiert die Gesellschaft:

\begin{table}[h]
\centering
\begin{tabularx}{\textwidth}{lXX}
\toprule
\textbf{Bereich} & \textbf{Heute} & \textbf{Mit systematischem Ansatz} \\
\midrule
Gesundheit & Präventionsprogramme werden wenig genutzt & Bessere Gestaltung erhöht Teilnahme \\
Umwelt & Appelle verhallen & Verhaltensbasierte Anreize wirken \\
Finanzen & Viele sparen zu wenig & Bessere Defaults helfen \\
Verwaltung & Angebote werden nicht genutzt & Bürgerfreundliche Gestaltung \\
\bottomrule
\end{tabularx}
\end{table}

\begin{summarybox}[Zusammenfassung: Impact for Switzerland]
\textbf{BEATRIX könnte der Schweizer Wirtschaft auf mehreren Ebenen nützen:}

\begin{itemize}[noitemsep]
    \item \textbf{Arbeitsplätze:} Konservativ geschätzt 30--50 direkt bis 2030
    \item \textbf{Know-how:} Schweizer Expertise in zukunftsrelevantem Feld
    \item \textbf{Modellcharakter:} Vorlage für KMU-Hochschul-Zusammenarbeit
    \item \textbf{Gesellschaft:} Bessere Entscheidungen in Gesundheit, Umwelt, Verwaltung
\end{itemize}

Die Schweiz hat eine einzigartige Ausgangslage durch die Kombination von Ernst Fehr, UZH und FehrAdvice. Dieses Potenzial verantwortungsvoll zu nutzen, ist das Ziel dieses Projekts.
\end{summarybox}

\newpage

%==============================================================================
% GESAMTZUSAMMENFASSUNG
%==============================================================================
\section*{Gesamtzusammenfassung}
\addcontentsline{toc}{section}{Gesamtzusammenfassung}

\begin{table}[h]
\centering
\begin{tabularx}{\textwidth}{lX}
\toprule
\textbf{Frage} & \textbf{Kernaussage} \\
\midrule
\textbf{a) Innovation} & BEATRIX macht verhaltensökonomische Analyse systematisch und nachvollziehbar -- durch Kontext-Modellierung, Wechselwirkungs-Analyse und transparente Entscheidungslogik \\
\textbf{b) Markt} & Die ``Soft-Factor-Lücke'' ist real: Für Finanzen gibt es SAP, für Verhalten fehlen systematische Werkzeuge \\
\textbf{c) Machbarkeit} & Das System funktioniert bereits -- der Innovation Cheque prüft Skalierbarkeit und Nutzerakzeptanz \\
\textbf{d) Partnerschaft} & UZH bringt genau die Kompetenzen, die FehrAdvice fehlen: technische Architektur und Akzeptanzforschung \\
\textbf{e) Schweiz} & Einzigartige Ausgangslage durch Fehr + UZH + FehrAdvice -- verantwortungsvoll nutzen \\
\bottomrule
\end{tabularx}
\end{table}

\vspace{1cm}

\begin{center}
\begin{tcolorbox}[colback=lightgray, colframe=innosuisse, width=0.95\textwidth]
\centering
\textbf{Der strategische Kernsatz}\\[0.5em]
\textit{``Dieses Projekt entwickelt keine neue Software, sondern eine wissenschaftlich fundierte Methodik für bessere Entscheidungen in Organisationen. Die Technologie ermöglicht diese Methodik -- aber die Transformation der Beratungspraxis ist das eigentliche Produkt.''}
\end{tcolorbox}
\end{center}

\vspace{2cm}

\begin{flushright}
\textit{Dokument erstellt: 18. Januar 2026}\\
\textit{Version: 2.0 -- Überarbeitet nach Style-Guide}\\
\textit{Für: Innosuisse Innovation Cheque (SSBM)}
\end{flushright}

\end{document}
